% ------------------------------------------------------------
% Page 98
% ------------------------------------------------------------

Du coin de l'âtre, une voix s'élève : « Et la guerre ? » Je comprends que ce seul mot
pose une question angoissante, et j'annonce la signature imminente des accords de
Munich ; j'assiste alors à une scène émouvante, que je ne puis éluder. Le père se lève,
embrasse ses deux fils, répétant : « Alors, io pas guéro ! ». La mère, les larmes aux
yeux, pousse un profond soupir.

Je suis en quelques instants celui qui apporte une bonne nouvelle. C'est dans une
atmosphère très détendue qu'on dresse la table sur laquelle on dispose tout ce qu'on a
de meilleur, et on débouche une bonne bouteille ; on a tout le temps maintenant de
faire connaissance... je ne rapporterai de cette très longue conversation que ce qui
m'intéressait directement.

« Vous arrivez de la Can de l'Hospitalet ! mais c'est tout proche ! En moins de trois
heures vous êtes chez vous. Justement, nous allons jeudi prochain à la foire de la saint
Michel à Barre (des Cévennes), et, si vous voulez venir avec nous, on vous montrera
le plus court chemin. Nous partirons à 5h00 et à 8 heures vous serez aux Crottes ! »
J'accepte avec empressement. J'apprends alors que les frères Rousset sont négociants
en bestiaux, répertoriés à Meyrueis, et non aux Oubrets. Ils se déplacent à pied pour
ramener leurs achats, d'ovins surtout.

La soirée se prolonge fort tard et je manifeste l'intention de me retirer à l'école.

-« L'école, ici on dit aussi le temple, c'est juste en face. Nous avons les clefs » ;

-« le temple ! ce mot me fait sursauter ;

-« un temple désaffecté sans doute ? »

-« Mais non, pas du tout ! le Pasteur de Meyrueis vient faire le culte deux ou trois
fois à la belle saison. La sacristie c'est votre logement ! ».

Ils perçoivent mon inquiétude, d'autant plus que je leur ai annoncé mon prochain
mariage.

-« Ecoutez, nous visiterons demain ; notre mère a renouvelé les draps et vous
pouvez aller vous reposer quand il vous plaira ».

Je m'étonne de tant de gentillesse.

-« Ne vous étonnez pas, c'est tous les ans que nous accueillons chez nous, parfois
pendant plusieurs semaines, les jeunes dans votre situation ; la chambre leur est
pratiquement réservée. L'an dernier c'était une demoiselle Melle. R..., de santé
délicate, souvent en congé de maladie. C'était trop dur pour elle ! ».

Effectivement, Melle. R..., originaire de l'Hérault me précédait d'un an à l'École
Normale, et je lui succédais un an plus tard aux Oubrets.

« Vous pouvez donc disposer de la chambre tant qu'il vous plaira, mais nous ne
pouvons pas vous nourrir. Nos heures ne coïncident pas avec les heures de l'école ;
Mais Mme. A..., qui ne peut pas vous loger, sera très heureuse de vous prendre en
pension ».
